\section{阶段性工作总结}
\label{sec:midterm-summary}

本章围绕多FPGA环境下的周期精确执行、分布式状态管理、自动化电路划分和运行控制，介绍了课题的阶段性研究进展。相较于仅将RTL设计部署到多个FPGA，本课题进一步关注如何在跨片执行后获得同一DUT逻辑周期的全局状态，并将这些状态用于恢复执行和软件RTL重放。目前已形成从层次化RTL处理、分区生成到双FPGA运行、检查点保存及恢复重放的处理流程，REMUS作为上述技术的实现与验证平台，已能够运行裸机程序和操作系统负载。

\paragraph{周期精确执行与分布式状态管理}

在DUOHUAFEN已有LI-TDM机制的基础上，REMUS优化了首尾传播阶段的收发流程，减少每个DUT逻辑周期中的冗余交换，并利用周期与传播阶段标识辅助检查通信序列。在此基础上，通过接收屏障、共同目标周期及差异化步进，实现各分区的一致暂停。该机制不要求主机请求同时到达，也不要求运行中的各分区始终具有相同进度，而是在提取状态前建立共同的逻辑时间边界。

沿用REMU已有的扫描、状态恢复和输入记录机制，已完成多分区扫描控制及检查点组织的扩展。各分区的局部扫描数据通过共同周期和编译生成的状态映射组成全局检查点，既可按分区写回FPGA，也可解码并加载到配套的软件RTL模型。由此，硬件执行、分布式状态提取与软件波形观测能够衔接起来。现有Cosim对照和板级恢复重放结果为这一流程提供了阶段性验证。

\paragraph{自动化划分与执行参数生成}

已实现依据RTL层次和资源规模选择性展开的自适应粒度处理，并形成混合粒度超图建模、求解和RTL分区生成流程。目前采用标量资源权重与KaHyPar完成划分，再利用组合依赖分析最大路径切割数，并据此配置跨分区传播阶段，将划分结果与周期执行联系起来。不同粒度的效率、分项资源的实际均衡程度及划分质量收益，仍需依据对照实验分别评价；标量表示导致的资源类型失衡也是后续需要改进的问题。

\paragraph{运行控制与Fabric扩展}

基于Host通过PCIe分别访问各FPGA的架构，已实现双FPGA系统的运行控制和检查点数据访问。针对主机直连接口数量对系统扩展的限制，进一步设计并实现了Fabric访问发起器、路由器、访问代理及Driver后端，建立由参考节点接入的统一资源访问结构和双节点Role联合仿真环境。该方向已有组件、软件访问及工程连接的阶段性验证，但完整联合访问验收和实际板卡验证仍需补充。

\paragraph{双FPGA系统集成与实验}

在两块基于VU37P的NM37板卡上，已通过PCIe主机访问和QSFP/Aurora板间连接完成原型系统集成。PicoRV32与Rocket能够正常执行所测试的裸机程序，香山南湖（Nanhu）与单核昆明湖（KMH）能够运行OpenSBI并进入Linux命令行。四个DUT均已完成一致暂停、检查点保存、FPGA恢复与重放以及软件RTL重放，说明前述控制与状态管理方法能够在不同规模设计上协同工作。

PicoRV32较长运行片段为0.326--0.327 MHz，KMH保存检查点前的运行片段为0.322 MHz。PicoRV32两条检查点保存记录约为84 ms，KMH的一次保存为80.1357 s，其中Host--DDR与文件处理阶段占约99.76\%。

划分方面已收集五个设计目录的15组2、3、4分区结果，求解调用耗时为0.136--8.122 s，最大路径切割数为2--5。另已统计45个Part的有效扫描状态位数和10个双分区Part的综合资源。综合结果显示，部分设计虽然LUT占用相近，BRAM、URAM或DSP仍集中在个别Part。因此，当前数据既支持已形成自动化处理与双板运行流程的阶段性结论，也为后续多维资源约束和通信性能分析提供了具体问题依据。

% TODO(系统材料)：Rocket/Nanhu若加入定量频率表，须补同口径日志；
% 当前总结只引用1.6已归档数字，不以历史单例填满四DUT矩阵。
% 平台和四个DUT的功能状态依据作者在1.6大纲讨论中的最新确认；
% 原始日志、配置、状态比较范围与重放波形仍需随1.6实验材料补齐。

总体而言，本阶段已将单FPGA检查点与重放能力扩展到双FPGA分区执行，并完成自动化划分、运行控制及板级系统验证所需的主要实现。下一阶段将在已有功能基础上扩大周期级状态比较和异常场景覆盖，校准资源估计并评价划分与通信开销，同时完善Fabric联合验证，逐步开展四FPGA等更大规模实验。
